\section{Construction and testing of pseudopotentials}\label{app:E.2}

What a pseudopotential gives up is the true valence wavefunction near the nucleus, exactly where the density vertex $\rho_{mn}(\vect{k}, \vect{q})$ of a heavy-mediator rate is evaluated, at momentum transfers of tens of keV. There a pseudopotential that is excellent for total energies, bands and phonons can fail without any sign in those quantities.

\subsection{The frozen core}

A pseudopotential acts on the valence states through a local radial potential $V_{\mathrm{loc}}(r)$ and a nonlocal part in the separable form of Kleinman and Bylander~\cite{kleinman1982efficacious},
\begin{equation}
  V_{\mathrm{NL}} = \sum_{\kappa}\sum_{ij}|\tilde p_{\kappa i}\rangle D_{\kappa,ij}\langle\tilde p_{\kappa j}| .
  \label{eq:E.2-KB}
\end{equation}
The projector functions $\tilde p_{\kappa i}(\vect{r})$, the nonlocal counterparts of the PAW projectors of \cref{sec:3.5}, are localized within a core radius $r_c$ of nucleus $\kappa$ and carry an angular momentum $l$ each. The coefficients $D_{\kappa,ij}$ form a small matrix. $V_{\mathrm{NL}}$ is the nonlocal pseudopotential of \cref{app:A.5}.

The approximation behind all of this is the frozen core. The core electrons are assumed to keep the orbitals they have in the isolated atom, unaffected by the crystal, the valence electrons, or a displacement of the nucleus. Its errors are of three kinds: transferability, the missing core region of the valence states, and the missing core states. Transferability means that a potential built for the atom must describe the atom in the crystal, at other bond lengths and charge states. Tests control this error. The missing core region of the valence states leaves total energies and bands almost untouched but changes matrix elements of operators that probe short distances, such as $\ee^{\ii\vect{q}\cdot\vect{r}}$ at large $q$. The all-electron reconstruction \cref{eq:3.5-rhoAE} repairs it. The core states themselves are absent from the spectrum, so a pseudopotential calculation of a response has no core edges. They are added back as atomic orbitals (\cref{app:E.7}). Core polarization by a displaced nucleus, which enters the large-$q$ lattice response, is identically zero in a frozen-core calculation and has to be computed separately (\cref{app:E.4}).

\subsection{Construction}

A pseudopotential is built from an all-electron calculation of the isolated atom, followed by the construction of pseudo-wavefunctions, the removal of the valence screening and the separable form. The atomic calculation is scalar-relativistic (that is, with the mass--velocity and Darwin terms of \cref{sec:3.5} but without spin--orbit coupling) and uses a reference configuration such as $[\mathrm{Ne}]\,3s^23p^2$ for silicon. The choice of valence comes with it. Silicon has an obvious split, four valence electrons above a core bound by about \SI{100}{eV} and more. Germanium does not: its 3d shell (\cref{sec:3.5}) overlaps the 4s and 4p orbitals in space and is polarized by the crystal.

The pseudo-wavefunction of each angular momentum $l$ equals the all-electron one beyond $r_c$, joins it smoothly at $r_c$, and is nodeless inside (\cref{sec:3.5}). Within these constraints, the shape inside $r_c$ is chosen to make the plane-wave expansion converge fast. In the optimized norm-conserving construction of Hamann~\cite{hamann2013optimized}, the pseudo-wavefunction is expanded in a few spherical Bessel functions inside $r_c$. The coefficients minimize the residual kinetic energy above a wavevector $q_{\mathrm{opt}}$~\cite[Eq.~(1)]{hamann2013optimized},
\begin{equation}
  E^{\mathrm{res}}(q_{\mathrm{opt}}) = \int_{q_{\mathrm{opt}}}^\infty\dd q\,q^4\,|\phi(q)|^2 , \qquad \phi(q) = 4\pi\int_0^\infty\dd r\,r^2j_l(qr)\,\phi(r) ,
  \label{eq:E.2-residual}
\end{equation}
subject to norm conservation and continuity at $r_c$. Here $\phi(r)$ is the radial pseudo-wavefunction and $j_l$ the spherical Bessel function of \cref{app:C.3}. Up to a constant, $q^4|\phi(q)|^2$ is the kinetic-energy density in momentum space, $q^2/2m_e$ times the radial measure $q^2|\phi(q)|^2$, so $E^{\mathrm{res}}$ is the kinetic energy missed by a basis truncated at $q_{\mathrm{opt}}$. The cutoff of the solid is then set by $q_{\mathrm{opt}}$, and the tail above it decays fast. The same construction uses two projectors per angular momentum, which reproduce the scattering of the all-electron atom at two energies, typically a bound state and a scattering state above it, rather than one.

Next, the valence part of the potential is removed. The potential that reproduces the pseudo-wavefunctions is a screened one: it contains the Hartree and exchange--correlation potentials of the valence density of the reference atom, which must be subtracted, since the crystal supplies its own. The Hartree potential is linear in the density and subtracts cleanly. The exchange--correlation potential is not, and where the core and valence densities overlap, subtracting the valence part alone leaves an error that changes with the valence configuration. The nonlinear core correction of Louie, Froyen and Cohen~\cite{louie1982nonlinear} keeps a smoothed partial core density in the exchange--correlation term for this reason. Finally, the semilocal potential is cast into the separable form \cref{eq:E.2-KB}, which reduces the cost of applying it from a double sum over plane waves to a few projections (\cref{eq:E.1-cost}).

Ultrasoft pseudopotentials and the PAW method relax these conditions to lower the cutoff. The ultrasoft scheme (\cref{sec:3.5}) pays for it with a generalized eigenvalue problem, and Kresse and Joubert showed how ultrasoft potentials follow from the PAW method as an approximation~\cite{kresse1999ultrasoft}. Norm-conserving potentials remain the simplest to use in density-functional perturbation theory, in electron--phonon calculations and in many-body perturbation theory, where every extra term of the ultrasoft and PAW formalisms has to be carried through each derivation~\cite{hamann2013optimized}.

The potentials of the representative calculations come from the SG15 library of Schlipf and Gygi~\cite{schlipf2015optimization}, generated with Hamann's construction, with parameters optimized automatically against all-electron results. The silicon potential has four valence electrons, s and p channels with $r_c = 2.40$ and $3.56$ bohr (\SI{1.27}{\angstrom} and \SI{1.88}{\angstrom}), two projectors per channel, $q_{\mathrm{opt}} = 5.11$ and $3.50$ inverse bohr, and no nonlinear core correction. The germanium potential has fourteen valence electrons, s, p and d channels with $r_c$ between 2.54 and 2.67 bohr, two projectors each, and $q_{\mathrm{opt}}$ up to 6.36 inverse bohr for the 3d channel. The squares of these wavevectors, \SI{26}{Ry} and \SI{12}{Ry} for silicon and \SI{40}{Ry} for the germanium 3d, are the scale at which the residual kinetic energy has been made small. They are below the cutoffs at which the total energy converges in \cref{tab:E-cutoff}, \SI{60}{Ry} and \SI{80}{Ry}, because the residual \cref{eq:E.2-residual} decays above $q_{\mathrm{opt}}$ but does not vanish.

\subsection{Testing}

A pseudopotential is tested on the atom, for transferability, and in the solid before it is used. The logarithmic derivative $\dd\ln R_l(r)/\dd r$ of the radial solution at a radius just beyond $r_c$, plotted against energy for the all-electron and the pseudo atom, measures how well the potential scatters at energies other than its reference energies. The two curves must agree over the whole window of energies the calculation will use, which for a response calculation includes empty states tens of eV above the gap. A separable potential can also bind a spurious state, a ghost, below or among the physical ones, because the eigenstates of \cref{eq:E.2-KB} are not ordered in energy by their number of nodes~\cite{hamann2013optimized}. A ghost appears as a spurious step in the logarithmic derivative, or in the crystal as a flat band. The PseudoDojo tables were scanned for ghosts up to about \SI{200}{eV} above the Fermi level, a range chosen for optical and $GW$ calculations, and a few heavy elements still show one in that range, between 20 and \SI{80}{eV}~\cite[Table~8]{vansetten2018pseudodojo}. Ghosts at still higher energies are unavoidable~\cite{vansetten2018pseudodojo}.

For transferability, the potential, built in one configuration, is applied to excited and ionized configurations of the atom, and the excitation energies are compared with all-electron values. In the solid, the most systematic test is the $\Delta$ gauge of Lejaeghere and coworkers. It is the root-mean-square difference between two equations of state $E(\Vol)$ of the same crystal~\cite[Eq.~(15)]{lejaeghere2014error},
\begin{equation}
  \Delta_{\mathrm{gauge}} = \biggl[\frac{1}{\Delta \Vol}\int_{\Vol_0 - \Delta \Vol/2}^{\Vol_0 + \Delta \Vol/2}\dd \Vol\,\bigl(E_1(\Vol) - E_2(\Vol)\bigr)^2\biggr]^{1/2} , \qquad \Delta \Vol = 0.12\,\Vol_0 ,
  \label{eq:E.2-delta}
\end{equation}
with both curves set to zero at their minima. The window spans $\pm\SI{6}{\percent}$ of the reference volume $\Vol_0$. It folds the errors of the volume, the bulk modulus and its pressure derivative into one number in meV per atom. Different all-electron codes agree with each other to 0.3--\SI{0.5}{meV} per atom on average over the periodic table. The norm-conserving PseudoDojo tables reach a mean of about \SI{0.7}{meV} per atom against an all-electron reference (excluding the magnetic Cr, Mn and Fe)~\cite{vansetten2018pseudodojo}. A potential must also be used with the functional with which it was generated; Hamann found that using a PBE-generated potential in an LDA calculation increased the differences from the all-electron results~\cite{hamann2013optimized}.

The standard tests are necessary but not sufficient for a dark-matter calculation. They probe the total energy and its derivatives with respect to volume, and the bands near the gap. Such a calculation also needs phonons, which probe the potential under small displacements of the nucleus. It needs empty states far above the gap, which probe the logarithmic derivatives at high energy. And it needs the density vertex at momentum transfers up to tens of keV, which probes the core region. The last of these is not part of any standard test.

\subsection{Semicore states}

When a shell sits between the deep core and the valence, its treatment changes both the accuracy and the cost. The germanium 3d shell is the case the book needs. Keeping it in the valence improves transferability, since the shell responds to its environment; it puts the 3d states into the band structure, where they contribute to the response above an energy deposit of about \SI{30}{eV} (\cref{sec:10.6}); and it costs. The fourteen electrons need seven occupied bands per atom instead of two, and the localized d states need a higher cutoff, \SI{80}{Ry} where silicon needs \SI{60}{Ry} (\cref{tab:E-cutoff}). Gallium arsenide, with the Ga 3d in the valence, is harder still: in the same kind of calculation the energy still changes by \SI{25}{meV} per atom between 100 and \SI{120}{Ry}. The alternative, a frozen 3d shell with a nonlinear core correction, gives a cheaper calculation of the valence and conduction bands and a response without the 3d channel, which must then be added as core states.

\subsection{The density vertex at large momentum}

The pseudo vertex is accurate for $qr_c \ll 1$ (\cref{sec:3.5}) and must be checked beyond. In the opposite limit, $q \gg q_{\mathrm{opt}}$, the pseudo-wavefunctions have almost no Fourier components by construction, \cref{eq:E.2-residual}. The pseudo vertex then falls off much faster than the true one, whose large-$q$ tail comes from the nodes and the cusp of the valence states near the nucleus. For the SG15 silicon potential $1/r_c$ is 1.6 and \SI{1.0}{keV}, and $q_{\mathrm{opt}}$ is 19 and \SI{13}{keV}. A dark-matter particle of \SI{10}{MeV} at the maximum speed in the Earth frame, $v_{\mathrm{max}} = 2.65\times10^{-3}$ (\cref{sec:1.3}), transfers at most $2m_\chi v_{\mathrm{max}} \approx \SI{53}{keV}$, and the typical momentum transfer of electron recoil is of the order of $\alpha m_e = \SI{3.7}{keV}$ (\cref{sec:1.3}). Much of the rate therefore comes from momentum transfers between $1/r_c$ and $q_{\mathrm{opt}}$, where the pseudo vertex of this potential is neither guaranteed accurate nor negligible, and a heavy-mediator rate weights the upper end of that range. A potential with a smaller $r_c$ pushes the problem to larger $q$ at the price of a higher cutoff.

In the reconstruction of Griffin and coworkers for silicon and germanium (\cref{we:3-Ecut}), the deeper core states, 2p and below in silicon and 3d and below in germanium in their treatment, were modeled as atomic orbitals in Bloch sums. The atomic description agreed with the reconstructed crystal states up to the reconstruction cutoff~\cite[Sect.~2.2]{griffin2021extended}.

\subsection{Nonlocal sum rules}

A nonlocal potential changes the commutators on which sum rules rest, and a test of a pseudopotential calculation against a sum rule must account for it. For a local potential, the double commutator of the Hamiltonian with $\ee^{\ii\vect{q}\cdot\vect{r}}$ involves only the kinetic energy and gives the $f$-sum rule of \cref{sec:4.3}. The plane wave $\ee^{\ii\vect{q}\cdot\vect{r}}$ is unitary. Inserting a complete set of states $|m\rangle$ of the one-electron Hamiltonian $H$ gives $\sum_m(\varepsilon_m - \varepsilon_n)|\langle m|\ee^{\ii\vect{q}\cdot\vect{r}}|n\rangle|^2 = \langle n|\ee^{-\ii\vect{q}\cdot\vect{r}}[H, \ee^{\ii\vect{q}\cdot\vect{r}}]|n\rangle = \langle n|\ee^{-\ii\vect{q}\cdot\vect{r}}H\ee^{\ii\vect{q}\cdot\vect{r}} - H|n\rangle$, and the same with $-\vect{q}$ in place of $\vect{q}$. Adding the two,
\begin{equation}
  \tfrac12\sum_m(\varepsilon_m - \varepsilon_n)\bigl(|\langle m|\ee^{\ii\vect{q}\cdot\vect{r}}|n\rangle|^2 + |\langle m|\ee^{-\ii\vect{q}\cdot\vect{r}}|n\rangle|^2\bigr) = \tfrac12\langle n|\ee^{-\ii\vect{q}\cdot\vect{r}}H\ee^{\ii\vect{q}\cdot\vect{r}} + \ee^{\ii\vect{q}\cdot\vect{r}}H\ee^{-\ii\vect{q}\cdot\vect{r}} - 2H|n\rangle .
  \label{eq:E.2-dc}
\end{equation}
The kinetic energy is shifted by the plane wave, $\ee^{-\ii\vect{q}\cdot\vect{r}}\vect{p}\,\ee^{\ii\vect{q}\cdot\vect{r}} = \vect{p} + \vect{q}$, so its contribution is $\tfrac12\langle(\vect{p} + \vect{q})^2 + (\vect{p} - \vect{q})^2 - 2\vect{p}^2\rangle/2m_e = q^2/2m_e$. A local potential commutes with the plane wave and drops out. A nonlocal one does not, and leaves
\begin{equation}
  \frac{q^2}{2m_e} + \delta_{\mathrm{NL}}(\vect{q}), \qquad
  \delta_{\mathrm{NL}}(\vect{q}) = \tfrac12\langle n|\ee^{-\ii\vect{q}\cdot\vect{r}}V_{\mathrm{NL}}\ee^{\ii\vect{q}\cdot\vect{r}} + \ee^{\ii\vect{q}\cdot\vect{r}}V_{\mathrm{NL}}\ee^{-\ii\vect{q}\cdot\vect{r}} - 2V_{\mathrm{NL}}|n\rangle ,
  \label{eq:E.2-fsumNL}
\end{equation}
on the right-hand side of \cref{eq:E.2-dc}. In the separable form \cref{eq:E.2-KB}, the conjugated potential $\ee^{-\ii\vect{q}\cdot\vect{r}}V_{\mathrm{NL}}\ee^{\ii\vect{q}\cdot\vect{r}}$ has the projectors $\ee^{-\ii\vect{q}\cdot\vect{r}}\tilde p_{\kappa i}$, so $\delta_{\mathrm{NL}}$ is a sum of projections and costs almost nothing to evaluate. At small $q$ it is of order $q^2$, like the kinetic term, and it does not vanish when all bands are summed: the pseudo-Hamiltonian obeys a modified $f$-sum rule. The all-electron Hamiltonian, whose potential is local, obeys the unmodified one. For a single Bloch state the same sum rule, without the symmetrization in $\pm\vect{q}$, keeps a term linear in $q$ set by the group velocity; \cref{eq:E.7-fsum} gives that form, and summed over a filled band it reduces to \cref{eq:E.2-dc}.

The representative silicon calculation shows the size of the effect. With the SG15 potential and the states at $\vect{k} = 0$, the left-hand side of \cref{eq:E.2-dc}, divided by $q^2/2m_e$, is 0.81 with 16 bands and 0.88 with 100 at $|\vect{q}| = \SI{0.1}{\angstrom^{-1}}$ along [100], and nearly the same at \SI{0.5}{\angstrom^{-1}}. The sum still rises slowly at 100 bands, so part of the deficit is band truncation, and part is $\delta_{\mathrm{NL}}$. The two must be separated, by computing \cref{eq:E.2-fsumNL} directly, before the sum rule can be used as a convergence test of the number of bands.

The limit $q \to 0$ of the same algebra gives the velocity, which carries the nonlocal commutator $\ii[V_{\mathrm{NL}}, \vect{r}]$ of \cref{app:A.5}; \cref{eq:E.7-vNL} gives its plane-wave form. The density vertex needs no such correction (\cref{sec:3.5}): a scalar potential couples to the density by multiplication, as in a local Hamiltonian~\cite[Eq.~(5)]{ismailbeigi2001coupling}.

The density vertex of a real crystal can now be computed at every momentum transfer that dark matter reaches. A rate, however, needs it on a fine grid of crystal momenta, where the bands near the gap vary on scales much finer than any affordable self-consistent calculation resolves. The localized orbitals of \cref{chap:25} turn a calculation on a coarse grid into one on any grid.
